\documentclass[11pt]{article}
\usepackage[T1]{fontenc}
\usepackage[utf8]{inputenc}
\usepackage[a4paper,margin=24mm]{geometry}
\usepackage{amsmath,amssymb,amsthm,mathtools,mathrsfs,booktabs,longtable,microtype,xurl}
\usepackage[hidelinks]{hyperref}
\newtheorem{theorem}{Theorem}[section]
\newtheorem{lemma}[theorem]{Lemma}
\newtheorem{proposition}[theorem]{Proposition}
\newtheorem{corollary}[theorem]{Corollary}
\newcommand{\C}{\mathbb C}
\newcommand{\Pp}{\mathcal P}
\newcommand{\Rr}{\mathcal R}
\newcommand{\tr}{\operatorname{Tr}}
\newcommand{\remm}{\operatorname{rem}}
\newcommand{\rank}{\operatorname{rank}}
\newcommand{\diag}{\operatorname{diag}}
\newcommand{\eps}{\varepsilon}
\newcommand{\Id}{I}
\newcommand{\norm}[1]{\left\lVert#1\right\rVert}
\setlength{\emergencystretch}{3em}
\title{The original observed arithmetic-current plane\\Complete conductor minima and prescribed signed classes}
\author{Zeta research continuation}
\date{22 September 2026}
\begin{document}
\maketitle
\begin{abstract}
The original rank-two arithmetic current is followed through the entire original observation minimum. A two-column lower bound retains the complex area of the two outgoing source classes; separate bounds for their norms would not suffice. On the fixed-period simple-quartet domain, both nonzero observed current eigenvalues have logarithmic magnitude $q\psi((N+1-q)/q)+O(k\log^2q)$, and have opposite signs. Two prescribed observed superpositions of the original source columns have those same signed logarithmic rates. The full kernel-current compression has norm at most $\exp[O(k\log^2q)]$, although the observed current has exponentially large positive and negative directions. All statements preserve the source cutoff, arithmetic unit, conductor moment and full quotient minima. The finite inequalities are proved here. The growing-degree rates use the retained original monic-profile theorem, with an explicit short extension of its degree range. They do not determine the current of the separately marked terminal eigenclass and do not prove RH.
\end{abstract}
\tableofcontents

\section{Original objects, source status, and the claim being proved}
Fix the stipulated simple quartet, the original admitted period $\varpi$ and branch, and the complete physical arithmetic unit. Keep
\[
 k\equiv1\pmod4,\quad k\ge17,\quad q=(k+1)^2,\quad q_0=(k-7)^2,
 \quad\Delta=q-q_0=16k-48,
\]
\[
 0<\delta<\tfrac12,\quad\gamma>2,\qquad
 Q_k(y)=\prod_{a,b=0}^k[y-(2b-k)\gamma+i(2a-k)\delta].
\]
The original coefficient algebra, arithmetic operator and physical coordinate are
\[
 E=\C[y]/Q_k,\qquad M[p]=[yp],\qquad
 S=k/2+iy,\qquad \mathscr A=kI/2+iM.
\]
The source is the unchanged even measure
\[
 d\mu_k(y)=w_h^{*k}(y)\,dy,\qquad
 w_h(y)=\frac{|(2\xi/h)(1/2+iy)|^2}{2\pi}.
\]
For $q-1\le N\le2q$, let $G_N$ be the entire degree-$N$ quotient minimum on $E$. The unchanged onto observation and its actual minimum section are
\[
 \Lambda:E\to B,\quad K=\ker\Lambda,\quad
 Q_{B,N}=(\Lambda G_N^{-1}\Lambda^*)^{-1},\qquad
 L_N=G_N^{-1}\Lambda^*Q_{B,N}.
 \tag{OCP1}
\]
The projection $P_{B,N}=L_N\Lambda$ is orthogonal in $G_N$. Its complement is $P_{K,N}$. Every adjoint below is either an explicitly written metric adjoint or the coordinate conjugate transpose $*$.

The reference measure and its complete monic norms are
\[
 d\sigma(y)=\frac{|\Gamma(1/4+iy/2)|^2}{2\pi}\,dy,\qquad
 h_n=\sqrt{2\pi}\,n!(1/2)_n.
 \tag{OCP2}
\]
Its monic polynomials $\pi_n$ satisfy
\[
 \pi_{n+1}=y\pi_n-n(n-1/2)\pi_{n-1},\qquad
 \sum_{n\ge0}\pi_n(y)t^n/n!=(1+t^2)^{-1/4}e^{y\arctan t}.
 \tag{OCP3}
\]
These are the Meixner--Pollaczek identities in the stated coordinate and mass; the ordinary identities are DLMF 18.22.8 and 18.23.7 \cite{dlmfpoly}. No monic norm or source measure is normalized to probability in the proofs.

Let $p_n$ be the real monic polynomial for $\mu_k$, with squared norm $\omega_n$. Use the complete comparison, through degree $2q+2$,
\[
 \ell\norm P_\sigma^2\le\norm P_{\mu_k}^2\le u\norm P_\sigma^2,
 \qquad \log(u/\ell)=O_h(k\log^2(q+2)).
 \tag{OCP4}
\]
The extra raising degree is used only to bound $p_{N+1}$ and its conductor image. The actual source metric remains the degree-$N$ minimum. The exact source identities and original monic-profile estimate are retained inputs \cite{heat,ivo}. The current's rank-two formula itself is rederived below; it is not claimed as a new identity. The quantitative new step is the full two-column inequality in Section 4.

\section{Two exact polynomial minima and the original current}
For the complete relation measure $|Q_k|^2d\mu_k$, let $\nu_j^\mu$ be the squared norm of its real monic polynomial $U_j$. The polynomial $Q_kU_j$ is monic of degree $q+j$ and orthogonal to all preceding original relation rows. Put
\[
 b_n=[p_n],\qquad E_N=\norm{b_N}_{G_N}^2,\qquad
 F_N=\norm{b_{N+1}}_{G_N}^2.
\]
\begin{lemma}
The complete original quotient norms are
\[
 E_N=\begin{cases}\omega_N,&N=q-1,\\
 \omega_N(1-\chi_N),&N\ge q,\end{cases}
 \qquad \chi_N=\omega_N/\nu_{N-q}^\mu,
\]
\[
 F_N=\nu_{N+1-q}^\mu-\omega_{N+1},\qquad
 \langle b_N,b_{N+1}\rangle_{G_N}=0.
 \tag{OCP5}
\]
The current is exactly
\[
 M=C_N+\frac{b_{N+1}b_N^*G_N}{\omega_N},\qquad
 C_N^{\dagger_{G_N}}=C_N.
 \tag{OCP6}
\]
The rank-one term squares to zero and has norm
\[
 \eps_N=\frac{\sqrt{E_NF_N}}{\omega_N}.
 \tag{OCP7}
\]
\end{lemma}
\begin{proof}
Projection of $p_N$ onto the complete relation space uses only its last orthogonal monic row: its coefficient is $\omega_N/\nu_{N-q}^\mu$. This proves the first formula, including the first cutoff where there is no relation row. The minimum representative of $b_{N+1}$ is
\[
 p_{N+1}-Q_kU_{N+1-q}\in\Pp_N.
\]
It is orthogonal to every relation of degree at most $N$. Its squared norm is $\nu_{N+1-q}^\mu-\omega_{N+1}$. The two classes have opposite parity in the even source and even relation quotient, proving orthogonality.

For the minimum representative $P$ of any class, project $yP$ onto $\Pp_N$. Its remaining polynomial is $[y^N]P\,p_{N+1}$. The further compression to the minimum-section space is selfadjoint, because multiplication by real $y$ is selfadjoint in the original source. Since
\[
 [y^N]P=\langle p_N,P\rangle_{\mu_k}/\omega_N
 =\langle b_N,[P]\rangle_{G_N}/\omega_N,
\]
this proves (OCP6). Orthogonality proves its square-zero assertion and (OCP7).
\end{proof}

Write $n_B=\dim B$ and define the actual observed centered current, in its original metric's orthonormal coordinates,
\[
 \mathscr W_{B,N}=Q_{B,N}^{-1/2}L_N^*
 (\mathscr A^*G_N+G_N\mathscr A-kG_N)
 L_NQ_{B,N}^{-1/2}.
 \tag{OCP8}
\]
Let
\[
 e_N=b_N/\sqrt{E_N},\qquad f_N=b_{N+1}/\sqrt{F_N},
 \quad U_N=[Q_{B,N}^{1/2}\Lambda e_N,\ Q_{B,N}^{1/2}\Lambda f_N].
\]
Its complete two-column Gram is
\[
 \mathcal C_N=U_N^*U_N=
 \begin{pmatrix}a_N&r_N\\\bar r_N&b_N^{\circ}\end{pmatrix},
 \qquad 0\preceq\mathcal C_N\preceq I_2.
 \tag{OCP9}
\]
The superscript on $b_N^{\circ}$ distinguishes the scalar norm from the quotient vector $b_N$. Define
\[
 J_2=\begin{pmatrix}0&-i\\i&0\end{pmatrix}.
\]
Then
\[
 \mathscr W_{B,N}=\eps_NU_NJ_2U_N^*.
 \tag{OCP10}
\]
Indeed $\mathscr A^*G+G\mathscr A-kG=i(GM-M^*G)$. This fixes the imaginary phase in the real-$y$ convention. In the physical monic convention $b_n^S=i^nb_n$, the corresponding rank-two matrix has the usual symmetric plus form.

\section{The original conductor and its two lower test classes}
Retain the complete original conductor
\[
 \mathcal Tp=\sum_{u,v=0}^8a_{uv}p(y+\sigma_{uv}),\qquad
 \sigma_{uv}=(2v-8)\gamma-i(2u-8)\delta.
\]
Its first nonzero moment has original order $v_0\le80$ and coefficient $\mu_{v_0}\ne0$. On a monic degree-$n$ polynomial its leading coefficient after lowering degree is
\[
 \kappa_n=i^{-v_0}\mu_{v_0}\binom n{v_0}.
 \tag{OCP11}
\]
Let $Q_0=Q_{k-8}$, $E_0=\C[y]/Q_0$, and let $C_A:E\to E_0$ be the induced map. The exact original identities are
\[
 C_A[p]=[\mathcal Tp]_{Q_0},\qquad K\subset\ker C_A.
 \tag{OCP12}
\]
No equality between these kernels is used. In particular no period-dependent choice of a generic kernel is made.

One finite forward bound from the full generating function is
\[
 \norm{\mathcal T}_{\sigma,\Pp_n\to\Pp_{n-v_0}}
 \le\mathcal F_n:=A_1e^{1+2\pi\gamma}(n+1)(2n+1)^{4\delta+1/2},
 \qquad A_1=\sum|a_{uv}|.
 \tag{OCP13}
\]
Translation multiplies (OCP3) by $e^{\sigma_{uv}\arctan t}$. On $|t|=n/(n+1)$, its modulus is at most $e^{2\pi\gamma}(2n+1)^{4\delta}$. Cauchy's estimate costs at most $e$, the full orthonormal coefficient ratio at most $\sqrt{2n+1}$, and summing all $n+1$ diagonals proves (OCP13). The original source version is the IVO forward bound \cite{ivo}.

For every class $x\in E$, minimize over its complete original observed fibre. Every candidate has degree at most $N$ and conductor value $C_Ax$. Consequently
\[
 \norm{\Lambda x}_{Q_{B,N}}^2
 \ge\frac{\ell}{\mathcal F_N^2}
\norm{C_Ax}_{G^-_D}^2,
 \quad D=N-v_0,
 \tag{OCP14}
\]
where $G^-_D$ is the full Gamma quotient metric on $E_0$ at degree $D$. All original upper relations map into the lower relation ideal, which is why (OCP14) holds for every representative.

Let $\nu_j^-$ denote the complete monic norm for $Q_0^2d\sigma$ and set
\[
 t=D-q_0,\qquad
 d_N=\frac{h_D}{\nu_t^-},\qquad
 e_N^{\circ}=\frac{h_{D+1}}{\nu_{t+1}^-}.
 \tag{OCP15}
\]
Here $D\ge q_0$ throughout the stated window when $k\ge17$. Both ratios are strictly between zero and one: equality would make an ordinary real-rooted orthogonal polynomial divisible by $Q_0$, whose roots are nonreal.

The lower quotient classes $a=[\pi_D]$, $b=[\pi_{D+1}]$ are orthogonal and have norms
\[
 A_D^2=h_D(1-d_N),\qquad
 B_D^2=h_{D+1}\frac{1-e_N^{\circ}}{e_N^{\circ}}.
 \tag{OCP16}
\]
For every $R\in\Pp_{D-1}$,
\[
 \left|\left\langle\frac a{A_D},[R]\right\rangle_{G^-_D}\right|
 \le\sqrt{\frac{d_N}{1-d_N}}\norm R_\sigma.
 \tag{OCP17}
\]
To prove this, the projection of $\pi_D$ onto the whole relation space is
$(h_D/\nu_t^-)Q_0U_t^-$. Its squared norm is $h_Dd_N$. The original $\pi_D$ is orthogonal to $R$; subtracting this exact relation projection and dividing by $A_D$ proves (OCP17).

\section{A finite area bound, retaining the complex cross term}
Write
\[
 \mathcal Tp_n=\kappa_n\pi_{n-v_0}+R_n,
 \qquad \deg R_n<n-v_0.
\]
By the full source comparison and orthogonality of the leading Gamma component,
\[
 \norm{R_n}_\sigma\le\norm{\mathcal Tp_n}_\sigma
 \le P_n:=\mathcal F_n\sqrt{u/\ell}\sqrt{h_n}.
 \tag{OCP18}
\]
Define the entirely specified positive scalar
\[
 \beta_n=\frac{P_n}{|\kappa_n|\sqrt{h_{n-v_0}}}.
\]
Let
\begin{align*}
 x_N&=\frac{\beta_N\sqrt{d_N}}{1-d_N},\\
 y_N&=\beta_{N+1}\sqrt{\frac{e_N^{\circ}}{1-e_N^{\circ}}},\\
 z_N&=\beta_N\beta_{N+1}
 \sqrt{\frac{e_N^{\circ}}{(1-d_N)(1-e_N^{\circ})}},\\
 \zeta_N&=(1-x_N)(1-y_N)-z_N.
\end{align*}
The symbol $\zeta_N$ here is a finite positive margin under the guard below, not a zeta zero or an observation pairing.

\begin{theorem}[Full observed two-column area]
Suppose
\[
 x_N+y_N+z_N\le\tfrac12.
 \tag{OCP19}
\]
Let $\nu_r^{+,\sigma}$ be the complete original $Q_k^2d\sigma$ monic norm, where $r=N+1-q$, and put
\[
 \xi_N=\frac{h_{N+1}}{\nu_r^{+,\sigma}},\qquad
 L_n=|\kappa_n|^2\frac{h_{n-v_0}}{h_n}
 =\frac{|\mu_{v_0}|^2}{(v_0!)^2}
 \prod_{j=0}^{v_0-1}\frac{n-j}{n-j-1/2}.
 \tag{OCP20}
\]
Then $\zeta_N\ge1/2$ and
\[
 \boxed{\Theta_N I_2\preceq\mathcal C_N\preceq I_2,\qquad
 \det\mathcal C_N\ge\Theta_N>0,}
 \tag{OCP21}
\]
where
\[
 \boxed{\Theta_N=\min\left\{1,
 \left(\frac\ell u\right)^2\mathcal F_N^{-4}
 L_NL_{N+1}(1-d_N)(1-e_N^{\circ})
 \frac{\xi_N}{e_N^{\circ}}\zeta_N^2\right\}.}
 \tag{OCP22}
\]
All minima in this formula are the complete monic minima of the stated measures.
\end{theorem}
\begin{proof}
Test the two conductor images against the orthonormal lower classes $a/A_D,b/B_D$, giving the complex matrix $Y$. By (OCP17)--(OCP18), its entries obey
\[
 |Y_{11}-\kappa_NA_D|\le\sqrt{d_N/(1-d_N)}P_N,
 \qquad |Y_{21}|\le P_N,
\]
\[
 |Y_{12}|\le P_{N+1},\qquad
 |Y_{22}-\kappa_{N+1}B_D|\le P_{N+1}.
\]
The off-diagonal leading term in $Y_{12}$ is zero by the exact parity orthogonality of $a$ and $b$. No phase of $R_N$ or of $\kappa_n$ is changed. The triangle inequality for the two products in the determinant gives
\[
 |\det Y|\ge
 |\kappa_N\kappa_{N+1}|A_DB_D
 \{(1-x_N)(1-y_N)-z_N\}.
 \tag{OCP23}
\]
The guard ensures the factors used in this lower bound are positive.

For $X=[b_N,b_{N+1}]$, the form inequality (OCP14), and then orthogonal projection onto the two lower test directions, give
\[
 X^*\Lambda^*Q_{B,N}\Lambda X
 \succeq\frac\ell{\mathcal F_N^2}(C_AX)^*G^-_D(C_AX)
 \succeq\frac\ell{\mathcal F_N^2}Y^*Y.
\]
Both are two-dimensional positive forms. Hence
\[
 \det\mathcal C_N\ge
 \left(\frac\ell{\mathcal F_N^2}\right)^2
 \frac{|\kappa_N\kappa_{N+1}|^2A_D^2B_D^2\zeta_N^2}{E_NF_N}.
\]
The full monic comparisons give $E_N\le u h_N$ and $F_N\le u\nu_r^{+,\sigma}$. Substitution and cancellation of the identical Gamma masses yield (OCP22). Since $\mathcal C_N\preceq I_2$, its smaller eigenvalue is at least its determinant, proving (OCP21).
\end{proof}

\section[The first source direction and the kernel current]{The first source direction and the whole kernel current}
The first normalized source direction is even better preserved than (OCP21) alone states. Put
\[
 \alpha_N=\norm{P_{K,N}e_N}_{G_N}^2.
\]
Every degree-$N$ minimum representative of a class in $K$ has conductor image divisible by $Q_0$. If its leading coefficient is $c$, that image has leading coefficient $\kappa_Nc$, so its full lower source norm is at least $|\kappa_Nc|\sqrt{\nu_t^-}$. Applying (OCP13)--(OCP14) to the leading-coefficient functional on $K$ proves
\[
 \boxed{\alpha_N\le
 \frac{\omega_N^2\mathcal F_N^2}
 {E_N\ell|\kappa_N|^2\nu_t^-}.}
 \tag{OCP24}
\]
This uses the identity $\langle b_N,x\rangle_{G_N}=\omega_N[y^N](\text{minimum representative of }x)$ on the entire kernel.

For an entirely Gamma-based upper certificate, assume additionally, when $N\ge q$, that
\[
 \frac u\ell\frac{h_N}{\nu_{N-q}^{+,\sigma}}\le\frac12.
 \tag{OCP25}
\]
At $N=q-1$ this extra guard is empty. Then (OCP5), the monic comparison and (OCP20) imply
\[
 \boxed{\alpha_N\le\mathfrak a_N:=
 \min\left\{1,\,2\frac u\ell\frac{\mathcal F_N^2}{L_N}d_N\right\}.}
 \tag{OCP26}
\]
The finite formula (OCP24) is valid without (OCP25).

The full normalized source pair is orthonormal, so
\[
 a_N=1-\alpha_N,\qquad
 |r_N|^2\le\alpha_N\norm{P_{K,N}f_N}^2\le\alpha_N.
 \tag{OCP27}
\]
The complete kernel-current compression obeys
\[
 \boxed{\norm{\mathscr W_{K,N}}\le2\eps_N\sqrt{\mathfrak a_N},\qquad
 |\tr\mathscr W_{B,N}|\le2\eps_N\sqrt{\mathfrak a_N}.}
 \tag{OCP28}
\]
For the first inequality write the original current as
$i\eps_N(f_Ne_N^*-e_Nf_N^*)$ in a source isometry and project both factors into $K$. For the second use (OCP27) and the exact trace in Section 7. The source total trace is zero, so the two diagonal traces are opposite. The full off-diagonal block $i\eps_N(f_Ke_B^*-e_Kf_B^*)$ is retained; no invariance of $K$ is inferred.

\section{Evaluate all growing-degree errors}
The sole asymptotic input beyond the original full-source comparison is the retained monic-profile theorem:
\[
 \log\frac{\nu_j^{+,\sigma}}{h_{q+j}}
 =2q\psi(j/q)+O_h(k\log^2(q+2)),\qquad0\le j\le q+1.
 \tag{OCP29}
\]
Its analogue for $Q_0$ holds with $q,k$ replaced by $q_0,k-8$. This is the reference version of the source's H4/V4 calculation; its actual arithmetic version follows by the same complete monic comparison \cite{heat,lrc}. It retains its written analytic proof status. The finite results preceding this section do not require (OCP29).

We explicitly extend the needed lower index by $O(k)$ instead of silently using the old degree range. For the lower relation measure $Q_0^2d\sigma$, let $a_j^-=\sqrt{\nu_j^-/\nu_{j-1}^-}$. The digamma series (DLMF 5.7.6, \cite{dlmfgamma}) and the nonreal root distance give
\[
 |(\log\sigma)'|\le2+\pi/2,\qquad
 |(\log(Q_0^2\sigma))'|\le B_0:=2q_0/\delta+2+\pi/2.
\]
For the first bound, in the series for $\Im\psi(1/4+ib)$ the initial term has absolute value at most two, and the remaining decreasing sum is bounded by $\pi/2$. Integration by parts against consecutive monic relation polynomials gives
\[
 j\nu_{j-1}^-\le B_0\sqrt{\nu_j^-\nu_{j-1}^-}.
\]
The trial $yU_{j-1}^-$ and the full Gamma multiplication bound give the other side:
\[
 \boxed{\frac j{B_0}\le a_j^-\le2(q_0+j).}
 \tag{OCP30}
\]
All boundary terms vanish by exponential decay. The Gamma recurrence (OCP3) proves $\norm{yP}_\sigma\le2(n+1)\norm P_\sigma$ for degree at most $n$.

For $j=q_0+O(k)$, a consecutive change in
$\log(\nu_j^-/h_{q_0+j})$ is therefore $O_h(\log q)$. At most $O(k)$ extra changes reach all indices $t,t+1$ used above, including the final $N=2q$ column. Since $\psi$ is smooth near one, (OCP29) extends to them with the same $O_h(k\log^2q)$ bound.

The exact reference profile may be written using its complementary elliptic ratio $r_s\in(0,1)$ and modulus $\kappa_s=\sqrt{1-r_s^2}$:
\[
 1+s=\frac{E(\kappa_s)}{r_sK(\kappa_s)},\qquad
 \psi(s)=s\log\frac{\kappa_s}{E(\kappa_s)}
 +\log\frac{1+r_s}{E(\kappa_s)}.
 \tag{OCP31}
\]
It is positive and decreasing on the required interval, with
\[
 a_0=\psi(0)=\log(4/\pi),\qquad
 a_1=\psi(1)=\log\frac{\kappa_1(1+r_1)}{E(\kappa_1)^2}>0.
\]
The endpoint expansion gives $|\psi'(s)|\le C(1+|\log s|)$ near zero; on compact positive intervals it is bounded. These are the same profile and elliptic conventions as the source, with the modulus rather than its square used in $K,E$ \cite{dlmfell,heat}.

Let $s_N=(N+1-q)/q$. The two indices of the lower profile differ by one, and
\[
 q_0\psi\left(\frac{r+\Delta-v_0}{q_0}\right)
 -q\psi(r/q)=O_h(k\log(q+2))
 \tag{OCP32}
\]
uniformly for $0\le r\le q+1$. To prove this, the two arguments differ by $O(k/q)$; integrate the logarithmic derivative bound over that interval. The factor $q-q_0=O(k)$ contributes only $O(k)$.

Equations (OCP18), (OCP20) and (OCP29)--(OCP32) yield, uniformly in the original window,
\[
 \log\beta_n=O_{h,\varpi}(k\log^2(q+2)),\qquad
 \log d_N,\log e_N^{\circ}
 =-2q\psi(s_N)+O_{h,\varpi}(k\log^2(q+2)).
\]
Consequently the finite guards (OCP19) and (OCP25) hold eventually at every cutoff. Crucially, in (OCP22) the ratio $\xi_N/e_N^{\circ}$ has logarithm
\[
 2q_0\psi\left(\frac{r+\Delta-v_0}{q_0}\right)
 -2q\psi(r/q)+O_h(k\log^2q)
 =O_{h,\varpi}(k\log^2q).
\]
Thus the two large monic factors have been compared before being bounded separately. We obtain
\[
 \boxed{0\le-\log\Theta_N=O_{h,\varpi}(k\log^2(q+2)),}
 \tag{OCP33}
\]
\[
 \boxed{\mathfrak a_N\le
 \exp[-2q\psi(s_N)+O_{h,\varpi}(k\log^2(q+2))].}
 \tag{OCP34}
\]
The inherited exact allowance formula (OCP5)--(OCP7) and (OCP29) also give
\[
 \log\eps_N=q\psi(s_N)+O_h(k\log^2(q+2)).
 \tag{OCP35}
\]
All implicit constants are fixed with the original packet and period; no uniformity near a vanishing conductor moment is claimed.

\section{Two observed signs and the complete kernel-current bound}
Let $r_N=x_N^{\circ}+iy_N^{\circ}$ in (OCP9). The nonzero spectrum of (OCP10) is that of $\eps_NJ_2\mathcal C_N$, whose characteristic polynomial is
\[
 \lambda^2+2\eps_Ny_N^{\circ}\lambda
 -\eps_N^2\det\mathcal C_N.
\]
The two eigenvalues are therefore exactly
\[
 \boxed{\lambda_{\pm,N}=\eps_N\left[
 -\Im r_N\pm\sqrt{a_Nb_N^{\circ}-(\Re r_N)^2}\right].}
 \tag{OCP36}
\]
They have opposite signs under (OCP19), and every remaining observed eigenvalue is zero. Orthogonal compression of the full current bounds their absolute values by $\eps_N$. Their product is $-\eps_N^2\det\mathcal C_N$, so
\[
 \boxed{\eps_N\Theta_N\le\lambda_{+,N},-\lambda_{-,N}\le\eps_N.}
 \tag{OCP37}
\]
The exact inertia is $(1,1,n_B-2)$, with all zero directions retained.

\begin{theorem}[Original observed current and hidden current]
For fixed admitted data and sufficiently large $k$, uniformly in $q-1\le N\le2q$,
\[
 \boxed{\begin{aligned}
 \log\lambda_{+,N}&=q\psi(s_N)+O_{h,\varpi}(k\log^2q),\\
 \log(-\lambda_{-,N})&=q\psi(s_N)+O_{h,\varpi}(k\log^2q).
 \end{aligned}}
 \tag{OCP38}
\]
In the entire original kernel,
\[
 \boxed{\norm{\mathscr W_{K,N}}+|\tr\mathscr W_{B,N}|
 \le\exp[O_{h,\varpi}(k\log^2(q+2))].}
 \tag{OCP39}
\]
The two observed signs are asymptotically balanced:
\[
 \left|\frac{\lambda_{+,N}}{-\lambda_{-,N}}-1\right|
 \le\frac{2\sqrt{\mathfrak a_N}}{\Theta_N}
 =\exp[-q\psi(s_N)+O_{h,\varpi}(k\log^2q)].
 \tag{OCP40}
\]
\end{theorem}
\begin{proof}
Equations (OCP33), (OCP35), and (OCP37) prove (OCP38). Equations (OCP28), (OCP34), and (OCP35) prove (OCP39). Finally $\lambda_++\lambda_-=-2\eps_N\Im r_N$ and $|r_N|\le\sqrt{\mathfrak a_N}$; divide by (OCP37) to prove (OCP40).
\end{proof}

This is the centered Hermitian current, not the full kernel action and not its energy. In particular, the previously proved growing rank of $\Lambda MI_K$ is compatible with (OCP39). The complete mixed current is
\[
 i\eps_N(f_Ke_B^*-e_Kf_B^*).
\]
Its difference from the first rank-one term has norm at most $\eps_N\sqrt{\mathfrak a_N}=e^{o(q)}$. This retains both leakage directions rather than assigning invariance to the kernel.

The latest PEN section isomorphisms equip the canonical section spaces with the actual pullbacks of restricted $G_N$. The current on their hidden domains is the restriction of $i(G_NM-M^*G_N)$. Equation (OCP39) therefore gives the same $e^{o(q)}$ operator-norm bound on every such hidden section domain. It does not replace the full ambient short minimum by the hidden minimum; PEN18--27 and its terminal correction remain unchanged \cite{pen}.

\section[Prescribed signed observed classes]{Two prescribed original observed classes with opposite current values}
The current signs can be attained without choosing its native eigenvectors. In the physical monic convention, define the cutoff-labelled original observed values
\[
 \mathfrak b_{\pm,N}
 =\Lambda\left(\frac{b_N^S}{\sqrt{E_N}}
 \pm\frac{b_{N+1}^S}{\sqrt{F_N}}\right),
 \qquad b_n^S=i^nb_n.
 \tag{OCP41}
\]
They differ by the common unit phase $i^N$ from the real-coordinate values $\Lambda(e_N\pm if_N)$. Their degree-$N$ source representatives use exactly the minimum polynomial of Section 2, including the cancellation of the degree-$N+1$ monic terms. No source cutoff is increased.

Let
\[
 \Phi_N(b)=b^*L_N^*(\mathscr A^*G_N+G_N\mathscr A-kG_N)L_Nb,
 \qquad j_{\pm,N}=\frac{\Phi_N(\mathfrak b_{\pm,N})}
 {\norm{\mathfrak b_{\pm,N}}_{Q_{B,N}}^2}.
\]
The complete complex formula is
\[
 \boxed{j_{\pm,N}=\eps_N\left[
 -2\Im r_N\ \pm\frac{2\det\mathcal C_N}
 {a_N+b_N^{\circ}\mp2\Im r_N}\right].}
 \tag{OCP42}
\]
It follows by substituting $(1,\pm i)^T$ into $\mathcal C_NJ_2\mathcal C_N$. In particular the cross phase is not replaced by the determinant.

If, in addition to (OCP19) and (OCP25),
\[
 \sqrt{\mathfrak a_N}\le\Theta_N/4,
 \tag{OCP43}
\]
then these two values are nonzero and
\[
 \boxed{\frac{\eps_N\Theta_N}{2}
 \le j_{+,N},-j_{-,N}\le\eps_N.}
 \tag{OCP44}
\]
To prove it, their squared observed norms are
$a_N+b_N^{\circ}\mp2\Im r_N$, lying in $(0,2]$ because $0\prec\mathcal C_N\preceq I$. Thus the positive fraction in (OCP42) is at least $\det\mathcal C_N$, whereas $2|\Im r_N|\le\Theta_N/2$. The upper bound is the full compressed operator norm bound.

Condition (OCP43) holds eventually by (OCP33)--(OCP34), uniformly over all original cutoffs. Therefore
\[
 \boxed{\log j_{+,N}=q\psi(s_N)+O_{h,\varpi}(k\log^2q),\qquad
 \log(-j_{-,N})=q\psi(s_N)+O_{h,\varpi}(k\log^2q).}
 \tag{OCP45}
\]
These are actual signed values on specified original observed classes. They are a cutoff-labelled source family, not the separately fixed terminal arithmetic eigenclass $v_\lambda$. Nothing in (OCP45) assigns a sign to $\Phi_N(\Lambda v_\lambda)$.

\section{Evaluated return and scalar allowance}
For the original four signs at $q-1,q,2q-1,2q$, (OCP38) and (OCP45) give
\[
 \boxed{\Rr\log\lambda_{+,N}
 =\Rr\log(-\lambda_{-,N})
 =2(a_0-a_1)q+O_{h,\varpi}(k\log^2q).}
 \tag{OCP46}
\]
The same formula holds with $\lambda_+$ and $-\lambda_-$ replaced by $j_+$ and $-j_-$. These are returns of cutoff-labelled operator spectra or cutoff-labelled classes, respectively. They are not a fixed-class determinant return.

At $N=q-1$, $s_N=0$; at $N=q$, $s_N=1/q$; and the upper two arguments are $1,1+1/q$. The endpoint derivative bound makes their small displacement cost only $O(\log q)$, already included in (OCP46).

A new outward integer calculation encloses
\[
\begin{split}
a_0&=0.2415644752704904446910368915632944245\ldots,\\
0.0665189520202739196651&<a_1<0.0665189520202739196824,
\end{split}
\]
and proves
\[
 \boxed{0.35009104650043305001
 <2(a_0-a_1)<0.35009104650043305006.}
 \tag{OCP47}
\]
The checker uses the actual root equation $E=2rK$, 4096 positive series terms with geometric tails, rational Machin bounds for $\pi$, integer square-root enclosures, and an outward logarithm series. This certifies a universal reference coefficient, not a hypothetical zero or native period.

For any scalar $b\ge0$, the inequality $\mathscr W_{B,N}+bI\succeq0$ is equivalent to $b\ge-\lambda_{-,N}$. Thus the least scalar positive correction is
\[
 \boxed{b_N^{\rm opt}=-\lambda_{-,N}
 =\exp[q\psi(s_N)+O_{h,\varpi}(k\log^2q)].}
 \tag{OCP48}
\]
Uniformly through the original window, it exceeds every fixed polynomial in $q,k$ eventually. In particular the original polynomial benchmark $\delta q(k+1)/2$ cannot dominate this canonical observed current. This is a conclusion about that stated metric and observation. A different tensor-section metric, a proper-source construction, or a geometric comparison is not declared equivalent to it.

\section{A new short-time law for the actual arithmetic propagator}
The current estimate also controls the actual first-order arithmetic evolution, rather than the heat of its positive energy. Define, in the original observed metric,
\[
 \mathscr U_{B,N}(t)=\Lambda e^{itM}L_N,
 \qquad \mathscr V_{B,N}(t)=e^{itM_{B,N}},
 \quad M_{B,N}=\Lambda ML_N.
 \tag{OCP49}
\]
These are different operators: the first measures the full original evolution, while the second evolves the actual minimum-section compression. They are not identified at finite time.

In the original observed isometry,
\[
 M_{B,N}=C_{B,N}+R_{B,N},\qquad
 C_{B,N}=J_B^\dagger C_NJ_B=C_{B,N}^\dagger,
 \quad R_{B,N}=\eps_N v_Nu_N^*,
\]
where $u_N,v_N$ are the two columns of $U_N$. Put
\[
 \eps_{B,N}=\norm{R_{B,N}}
 =\eps_N\sqrt{a_Nb_N^{\circ}},\qquad
 z_{B,N}=\tr R_{B,N}=\eps_Nr_N.
 \tag{OCP50}
\]
Thus $R_{B,N}^2=z_{B,N}R_{B,N}$, and the previous finite inequalities give
\[
 \eps_N\sqrt{\Theta_N}\le\eps_{B,N}\le\eps_N,
 \qquad |z_{B,N}|\le\eps_N\sqrt{\mathfrak a_N}.
 \tag{OCP51}
\]
The full polynomial multiplication inequality gives the finite bound
\[
 \norm{C_N},\norm{C_{B,N}}\le
 \mathfrak c_N:=2(N+1)\sqrt{u/\ell},\qquad
 \log\mathfrak c_N=O_{h,\varpi}(k\log^2q).
 \tag{OCP52}
\]
It follows by applying the original source comparison to $yP$ and $P$, and the exact Gamma recurrence to $P$ through degree $N$. The source cutoff in (OCP49) is still $N$.

\begin{lemma}[Rank-one propagator estimate with the noncommuting part retained]
Let $C$ be any bounded operator, let $R$ have rank one and norm $\epsilon$, and write $R^2=zR$. For $T\ge0$ set
\[
 \delta_T=\norm C\,e^{T|z|}(T+\epsilon T^2/6).
\]
If $\delta_T<1$, then for $0\le t\le T$,
\[
 \boxed{
 \frac{\norm{e^{it(C+R)}-I-itR}}{1+t\epsilon}
 \le e^{T|z|}\left[
 \frac{\delta_T}{1-\delta_T}+\frac{T|z|}{2}\right].}
 \tag{OCP53}
\]
No commutation of $C$ and $R$ is assumed.
\end{lemma}
\begin{proof}
Exactly,
\[
 e^{itR}=I+f_t(z)R,\quad
 f_t(z)=\frac{e^{itz}-1}{z},\quad f_t(0)=it,
\]
so $\norm{e^{itR}}\le e^{T|z|}(1+t\epsilon)$ and
$|f_t(z)-it|\le t^2|z|e^{T|z|}/2$.
Duhamel's identity gives
\[
 e^{it(C+R)}-e^{itR}
 =i\int_0^t e^{i(t-s)R}C e^{is(C+R)}ds.
\]
Let $M_T=\sup_{0\le s\le T}\norm{e^{is(C+R)}}/(1+s\epsilon)$. The exact integral of the norm-ratio factor is
\[
 \int_0^t\frac{(1+(t-s)\epsilon)(1+s\epsilon)}{1+t\epsilon}ds
 =t+\frac{\epsilon^2t^3}{6(1+t\epsilon)}
 \le T+\epsilon T^2/6.
\]
Hence $M_T\le e^{T|z|}+\delta_TM_T$ and the relative Duhamel error is at most $\delta_Te^{T|z|}/(1-\delta_T)$. Adding the exact rank-one exponential remainder proves (OCP53).
\end{proof}

Use (OCP53) first on the full source, where $z=0$, and then on the observed compression, where $z=z_{B,N}$. The first error, after observation, must be multiplied by
\[
 \frac{1+t\eps_N}{1+t\eps_{B,N}}
 \le\Theta_N^{-1/2};
 \tag{OCP54}
\]
this factor is retained. The main term in both cases is exactly $I+itR_{B,N}$.

For every fixed $b>a_0/2$, take the same real time at all cutoffs,
\[
 t_k(b)=e^{-bq}.
\]
Equations (OCP33)--(OCP35), (OCP51)--(OCP54) give, for a fixed $c_b>0$ smaller than $\min\{b,2b-a_0\}$,
\[
 \boxed{
 \sup_N\frac{
 \norm{\mathscr U_{B,N}(t_k)-I-it_kR_{B,N}}
 +\norm{\mathscr V_{B,N}(t_k)-I-it_kR_{B,N}}
 }{1+t_k\eps_{B,N}}
 \le e^{-c_bq+O_{h,\varpi,b}(k\log^2q)}.}
 \tag{OCP55}
\]
In particular their difference has the same relative bound. This comparison includes the actual hidden-source coupling through Duhamel's complete integral; it is not a claim that observation commutes with evolution.

A rank-one operator $R/\epsilon$ of norm one and trace $\zeta$ is unitarily equivalent on its nontrivial two-dimensional span to
\[
 \begin{pmatrix}\zeta&\sqrt{1-|\zeta|^2}\\0&0\end{pmatrix}.
\]
Its distance from the unit nilpotent matrix is at most $2|\zeta|$. Here $|z_{B,N}|/\eps_{B,N}\le\sqrt{\mathfrak a_N/\Theta_N}$ is exponentially small. The nilpotent shear has exact norm
\[
 \left\|I+ix\begin{pmatrix}0&1\\0&0\end{pmatrix}\right\|
 =\frac{\sqrt{x^2+4}+x}{2},\qquad x\ge0.
 \tag{OCP56}
\]
Combining it with (OCP55) proves
\[
 \boxed{
 \norm{\mathscr U_{B,N}(t_k(b))}
 =\norm{\mathscr V_{B,N}(t_k(b))}(1+o(1))
 =\frac{\sqrt{4+x_{k,N}^2}+x_{k,N}}{2}(1+o(1)),
 \quad x_{k,N}=t_k(b)\eps_{B,N}.}
 \tag{OCP57}
\]
Both relative errors are uniform over the original cutoff window. Their rates follow from (OCP55) and the explicit $\sqrt{\mathfrak a_N/\Theta_N}$ error. The equations do not require an absolute evaluation of the finite $\eps_{B,N}$ to obtain the following logarithmic value:
\[
 \boxed{\begin{gathered}
 \frac1q\log\norm{\mathscr U_{B,N}(e^{-bq})}
 =\frac1q\log\norm{e^{ie^{-bq}M_{B,N}}}+o(1),\\
 \frac1q\log\norm{\mathscr U_{B,N}(e^{-bq})}
 =(\psi(s_N)-b)_++o_{h,\varpi,b}(1),\qquad b>a_0/2.
 \end{gathered}}
 \tag{OCP58}
\]
It includes equality $\psi(s_N)=b$, because
$\log[(\sqrt{4+x^2}+x)/2]$ differs from $\log(1+x)$ by at most a fixed constant.

In particular $a_1<1/8<a_0$ and $1/8>a_0/2$. At this completely specified common time,
\[
 \boxed{\begin{aligned}
 \Rr\log\norm{\Lambda e^{ie^{-q/8}M}L_N}_{Q_{B,N}}
 &=(2\log(4/\pi)-1/4)q+O_{h,\varpi}(k\log^2q),\\
 \Rr\log\norm{e^{ie^{-q/8}M_{B,N}}}_{Q_{B,N}}
 &=(2\log(4/\pi)-1/4)q+O_{h,\varpi}(k\log^2q).
 \end{aligned}}
 \tag{OCP59}
\]
The coefficient is
\[
 2\log(4/\pi)-1/4
 =0.2331289505409808893820737831265888490\ldots.
\]

The original uncentered physical evolution is $e^{t\mathscr A}=e^{kt/2}e^{itM}$. The common scalar adds $kt/2$ to each logarithmic norm and cancels exactly under the four signs. Thus (OCP59) is also the return of the actual physical scaling evolution, with its central factor retained.

This growth is not eigenvalue growth of the original arithmetic spectrum: every eigenvalue of $e^{it_kM}$ tends to one since $|\omega_{ab}|=O_h(k)$. The intrinsic observed generator has
\[
 \norm{M_{B,N}^2}\le
 \mathfrak c_N^2+2\mathfrak c_N\eps_{B,N}
 +|z_{B,N}|\eps_{B,N}\le e^{a_0q+o(q)}.
\]
Hence its evolution eigenvalues also tend uniformly to one when $b>a_0/2$, while its operator norm has the positive exponential rate in (OCP58) where $\psi(s_N)>b$. The original metric nonnormality, rather than an inserted spectrum, produces the evaluated transient.

\section{Verification and the precise remaining boundary}
The new checker reconstructs complete finite Gamma-polynomial source Grams, entire relation minima, nontrivial conductor maps, complex observations, and original quotient projections. Its auxiliary fixtures include conductor orders zero and one, output ranks two and three, nonunitary changes of coefficient frame, and arbitrary retained common source masses. Two auxiliary large-gap examples satisfy the finite determinant-dominance guard and check the full numerical value of the rational lower certificate. The native-shaped small fixtures test the finite identities without asserting the eventual guard there. Separate-column positivity and deletion of the real-coordinate imaginary phase are rejected controls.

The finite checks are not the analytic proof. The asymptotic theorem uses the original source comparison and monic-profile theorem. The latter is neither inferred from a heat determinant nor promoted from a finite auxiliary grid. The bounded additional lower-degree range is proved in (OCP30).

The prior responses already recorded formal rank-two spectra in unknown projected components and asserted related exponential behavior. This paper supplies an independent full-source two-column determinant proof, its explicit positivity guard, the kernel-current bound, and prescribed current-bearing observed classes. It does not present the inherited rank-two characteristic formula as a new discovery.

The latest consulted publication is commit \nolinkurl{2b1445b21cc2fcd7e8b7ff34c6aea7139fb6afb5}, edition SZ-20260922-024. Its kernel determinant, uniform measured-resolvent and measured frequency-phase results remain unchanged. Those frequency phases are not the first-order arithmetic current calculated here. The separate marked terminal-eigenclass current and the independent proper-source target minimum remain separate calculations; neither is assigned a value by the operator signature.

\begin{thebibliography}{9}
\bibitem{dlmfpoly} Koornwinder, T. H., Wong, R., Koekoek, R., \& Swarttouw, R. F. (n.d.). Orthogonal polynomials. In \emph{NIST Digital Library of Mathematical Functions}, Chapter 18, equations 18.22.8 and 18.23.7. \url{https://dlmf.nist.gov/18.22.E8}; \url{https://dlmf.nist.gov/18.23.E7}.
\bibitem{dlmfgamma} Askey, R. A., \& Roy, R. (n.d.). Gamma function. In \emph{NIST Digital Library of Mathematical Functions}, Chapter 5, equation 5.7.6. \url{https://dlmf.nist.gov/5.7}.
\bibitem{dlmfell} Carlson, B. C. (n.d.). Elliptic integrals. In \emph{NIST Digital Library of Mathematical Functions}, Chapter 19, section 19.4. \url{https://dlmf.nist.gov/19.4}.
\bibitem{heat} Zeta research continuation. (2026). \emph{The original exponential-time heat transition and the cost of weak source activation}, H1--H9. Complete supplied file \nolinkurl{COMPLETE_PROOFS (4).md}. The original $y$ convention, complete monic norms and retained monic-profile theorem are the inputs cited here.
\bibitem{ivo} Zeta research continuation. (2026). \emph{The original observation preserves an entire inverse-power determinant at the $kq$ scale}, equations 8--13. Complete supplied file \nolinkurl{Pasted markdown(20260921-031454).md}. The exact conductor, first nonzero moment and forward norm are retained.
\bibitem{lrc} Zeta research continuation. (2026). \emph{Displaced lower-root covariance and source-order continuation}. Complete supplied file \nolinkurl{Pasted markdown(20260919-191918).md}, monic estimates in section 1. The lower source order remains one.
\bibitem{pen} Split-Zero research programme. (2026). \emph{The original kernel pencil, canonical section maps, and the coupled physical current}, PEN14--PEN27. Commit \nolinkurl{9104d852c6bc5b79e6fc4588c403128a40b7fda3}, file \nolinkurl{workbenches/splitzero-tandem/continuations/20260922-arithmetic-probe-mass/PENCIL_PROOFS.md}.
\bibitem{current} Split-Zero research programme. (2026). \emph{Original kernel coefficient and measured frequency return}, SZ-20260922-024. Commit \nolinkurl{2b1445b21cc2fcd7e8b7ff34c6aea7139fb6afb5}, file \nolinkurl{workbenches/splitzero-tandem/continuations/20260922-kernel-frequency/RESULTS_20260922_024.md}.
\end{thebibliography}
\end{document}
